\chapter{Presupuesto económico}
\label{cap:presupuesto}

Este anexo estima el coste de realizar el trabajo. Se estructura en costes de personal,
de material y de licencias. Las cifras de personal y de amortización se cerrarán al final del
desarrollo, cuando el cómputo de horas dedicadas sea completo.

\section{Costes de personal}

El coste de personal se calculará como el producto de las horas dedicadas a cada fase por un
coste horario de referencia para un ingeniero recién titulado. Las horas se obtienen del
registro de trabajo del proyecto, que anota cada sesión con su objetivo y su resultado.

% TODO: tabla de horas por fase y coste horario de referencia.

\section{Costes de material}

Todo el desarrollo se ha realizado sobre una única estación de trabajo propia, con procesador
de 16 núcleos, 64~GB de memoria y GPU dedicada, sin infraestructura de nube ni de centro de
datos. Su coste se imputará como amortización lineal durante los meses de uso, sobre una vida
útil de referencia.

% TODO: precio de adquisición, vida útil y meses de uso.

\section{Costes de licencias}

El coste de licencias es \textbf{nulo}. Todos los componentes empleados son de código abierto
o de uso gratuito (Anexo~\ref{anx:licencias}), y los conjuntos de datos y las bases de
conocimiento son públicos. Las plataformas comerciales que ofrecen capacidades equivalentes,
como los SIEM con correlación y priorización de incidentes, se licencian habitualmente por
volumen de datos ingeridos o por equipo supervisado, lo que constituye precisamente la barrera
de entrada que el trabajo pretende reducir.

% \section{Resumen del presupuesto}
% TODO: tabla resumen (personal + material + licencias, costes generales, beneficio e IVA).
